\documentclass{article}
\usepackage[english]{babel}
\usepackage{graphicx}
\usepackage{tabulary}
\usepackage{tabularx}
\usepackage{fixltx2e}
\usepackage{framed}
\usepackage{tikz}
\usepackage{mwe}
\usepackage{xcolor}
%\textsubscript{this}
\usepackage[margin=1in,top=1.2in,headheight=57pt,headsep=0.1in]
{geometry}
\usepackage{fancyhdr}
\everymath{\displaystyle}
\usetikzlibrary{calc}
\usetikzlibrary{shapes.multipart, shapes.geometric, arrows}
\usepackage{soul}
\usepackage{cancel}
\usepackage[utf8]{inputenc}
%\usepackage[dvipsnames]{xcolor}
%\usepackage{xcolor}
%\textsubscript{this}
\usepackage{lastpage}
\usepackage[utf8]{inputenc}
\usepackage{ifthen}
\usepackage{amsmath}
\usepackage{fancyhdr}
\usepackage[document]{ragged2e}
\usepackage[margin=1in,top=1.2in,headheight=57pt,headsep=0.1in]
{geometry}
\usepackage{fancyhdr}
\everymath{\displaystyle}
\linespread{2}%controls the spacing between lines. Bigger fractions means crowded lines%
%\pagestyle{fancy}
%\usepackage[margin=1 in, top=1in, includefoot]{geometry}
%\everymath{\displaystyle}
\linespread{2}%controls the spacing between lines. Bigger fractions means crowded lines%
%\pagestyle{fancy}
\pagestyle{fancy}
\setlength{\headheight}{56.2pt}
\colorlet{Mycolor1}{green!10!orange!90!}
\definecolor{shadecolor}{RGB}{80,180,180}


\begin{document}
\chead{\ifthenelse{\value{page}=1}{\textbf Removal Efficiencies}{}}
\rhead{\ifthenelse{\value{page}=1}{Shabbir Basrai}{Shabbir Basrai}}
\lhead{\ifthenelse{\value{page}=1}}
\cfoot{}
\cfoot{\thepage\ of \pageref{LastPage}}
\renewcommand{\headrulewidth}{2pt}
\renewcommand{\footrulewidth}{1pt}
%Defining colour with different models.
%\definecolor{mypink1}{rgb}{0.858, 0.188, 0.478}
%\definecolor{mypink2}{RGB}{219, 48, 122}
%\definecolor{mypink3}{cmyk}{0, 0.7808, 0.4429, 0.1412}
%\definecolor{mygray}{gray}{0.6}
%\colorlet{LightRubineRed}{RubineRed!70!}
%\colorlet{Mycolor1}{green!10!orange!90!}
%\definecolor{Mycolor2}{HTML}{00F9DE}

%New command used in the table with all available colour names
\newcommand{\thiscolor}[1]{\texttt{#1} \hfill \fcolorbox{black}{#1}{\hspace{2mm}}}

%This changes the row separation in the table
\renewcommand{\arraystretch}{1.5}

%This document present several examples on how to use the \texttt{color} package to change the colour of elements in \LaTeX.

%\begin{itemize}
%\item \textcolor{Mycolor1}{First item}
%\item \textcolor{Mycolor2}{Second item}
%\end{itemize}

%\noindent
%{\color{LightRubineRed} \rule{\linewidth}{1mm} }

%\noindent
%{\color{RubineRed} \rule{\linewidth}{1mm} }


%Not only blocks, such as environments, can be set to a determined colour, but some \textcolor{red}{special words} too. You can even use your own user-defined colours. Below the same colour with different models:

%\begin{enumerate}
%\item \textcolor{mypink1}{Pink with rgb}
%\item \textcolor{mypink2}{Pink with RGB}
%\item \textcolor{mypink3}{Pink with cmyk}
%\item \textcolor{mygray}{Gray with gray}
%\end{enumerate}

%The background colour of some text can also be \textcolor{red}{easily} set. For instance, you can change to orange the background of \colorbox{BurntOrange}{this text} and then continue typing.

%\clearpage

%\begin{center}
%\begin{tabular}{ llll } 
%\thiscolor{Apricot} & \thiscolor{Emerald} & %\thiscolor{OliveGreen} & \thiscolor{RubineRed} \\ 
%\thiscolor{Aquamarine} & \thiscolor{ForestGreen} & \thiscolor{Orange} & \thiscolor{Salmon} \\ 
%\thiscolor{Bittersweet} & \thiscolor{Fuchsia} & %\thiscolor{OrangeRed} & \thiscolor{SeaGreen}\\ 
%\thiscolor{Black} & \thiscolor{Goldenrod} & %\thiscolor{Orchid} & \thiscolor{Sepia}\\ 
%\thiscolor{Blue} & \thiscolor{Gray} & %\thiscolor{Peach} & \thiscolor{YellowOrange}\\ 
%\thiscolor{BlueGreen} & \thiscolor{Green} & %\thiscolor{Periwinkle} & \thiscolor{SkyBlue}\\ 
%\thiscolor{BlueViolet} & \thiscolor{GreenYellow} & \thiscolor{PineGreen} & \thiscolor{SpringGreen}\\ 
%\thiscolor{BrickRed} & \thiscolor{JungleGreen} & \thiscolor{Plum} & \thiscolor{Tan}\\ 
%\thiscolor{Brown} & \thiscolor{Lavender} & \thiscolor{ProcessBlue} & \thiscolor{TealBlue}\\ 
%\thiscolor{BurntOrange} & \thiscolor{LimeGreen} & \thiscolor{Purple} & \thiscolor{Thistle}\\ 
%\thiscolor{CadetBlue} & \thiscolor{Magenta} & \thiscolor{RawSienna} & \thiscolor{Turquoise}\\ 
%\thiscolor{CarnationPink} & \thiscolor{Mahogany} & \thiscolor{Red} & \thiscolor{Violet}\\ 
%\thiscolor{Cerulean} & \thiscolor{Maroon} & \thiscolor{RedOrange} & \thiscolor{VioletRed}\\ 
%\thiscolor{CornflowerBlue} & \thiscolor{Melon} & \thiscolor{RedViolet} & \thiscolor{White}\\ 
%\thiscolor{Cyan} & \thiscolor{MidnightBlue} & \thiscolor{Rhodamine} & \thiscolor{WildStrawberry}\\ 
%\thiscolor{Dandelion} & \thiscolor{Mulberry} & \thiscolor{RoyalBlue} & \thiscolor{Yellow}\\ 
%\thiscolor{DarkOrchid} & \thiscolor{NavyBlue} & \thiscolor{RoyalPurple} & \thiscolor{YellowGreen}
%\end{tabular}%
%\end{center}%

\textbf{PROCESS REMOVAL RATES OR EFFICIENCY}\\
\begin{itemize}
\item Process removal rate or removal efficiency is the percentage of the inlet concentration removed.  
\item It is used for quantifying the pollutant removal during wastewater treatment and is established based upon the amount of a particular wastewater constituent entering and leaving a treatment process.

\item $Process \enspace Removal \enspace Rate \enspace (\%) = \frac{Pollutant \enspace  In-Pollutant\enspace  Out}{Pollutant \enspace In}*100$\\

\item If 10 units of a pollutant are entering a process and 8 units of pollutant are leaving (process removes 2 units), then the process removal rate for that pollutant is (10-8)/10*100=20\%.  In this example the process is 20\% efficient in removing that particular pollutant.

\item The amount of pollutant can be measured in terms of concentration (mg/l) or in terms of mass loading (lbs).  The pounds formula is used for calculating the mass loadings.  
\end{itemize}
The above example is for calculating the removal efficiency using the inlet and outlet concentrations or mass loading.\\
The methods below can be used for calculating either the inlet or outlet pollutant concentrations, if the removal efficiency and the corresponding inlet or outlet concentrations are given. 

\begin{snugshade*}
\noindent\textsc{Case 1:  Calculating outlet conc. (X) given the inlet conc. and removal efficiency (RE\%):}
\end{snugshade*}
\tikzstyle{block} = [rectangle, draw, fill=red!40, 
    text width=6em, text centered, rounded corners, minimum height=3em]
\tikzstyle{arrow} = [draw, -latex']
\begin{figure}[!h]
\centering
\begin{tikzpicture}[node distance =1.5cm, auto]
    \draw ++(0,0) node [block] (Process) {Process};
   \node[node distance=1.9in] (dummy_in) [left of=Process] {In};
   \node[node distance=1.9in] (dummy_out) [right of=Process] {Out};
	\node (Removal) [below of=Process, yshift=-0in] {$\tiny{Removal \enspace Efficiency=RE\% \enspace (Given)}$};
    \path [arrow] (dummy_in)-- (Process)  node [above] {\hspace{-5.8cm}$A \enspace mg/l \enspace (Given) $} node [below] {\hspace{-5.8cm}$100 \enspace mg/l$};
    \path [arrow] (Process) -- (dummy_out)  node [above] {\hspace{-4cm}$X \enspace mg/l \enspace (Unknown)$} node [below] {\hspace{-3.9cm}($100-RE\%)\enspace mg/l$};
   \draw[arrow] (Process) -- (Removal);
\end{tikzpicture}
\end{figure}
Using the fact that if the inlet concentration was 100 mg/l, the outlet concentration would be 100 minus the removal efficiency.\\
Setup the equation as:  $\frac{Out}{In}: \enspace \frac{X \enspace mg/l}{A \enspace mg/l}=\frac{100-RE\%}{100}$\\
\vspace{0.3cm}
Calculate X using cross multiplication - if $\frac{A}{B}=\frac{C}{D} \implies A=B*\frac{C}{D}$:\\
$X \enspace mg/l=A \enspace mg/l*\frac{100-RE\%}{100}$\\

 \vspace{3cm}
\begin{snugshade*}
\noindent\textsc{Case 2:  Calculating inlet conc. (X) given the outlet conc. and removal efficiency (RE\%):}
\end{snugshade*}





\begin{figure}[!h]
\centering
\begin{tikzpicture}[node distance =1.5cm, auto]
    \draw ++(0,0) node [block] (Process) {Process};
   \node[node distance=1.9in] (dummy_in) [left of=Process] {In};
   \node[node distance=1.9in] (dummy_out) [right of=Process] {Out};
	\node (Removal) [below of=Process, yshift=-0in] {$Removal \enspace Efficiency=RE\% \enspace (Given)$};
    \path [arrow] (dummy_in)-- (Process)  node [above] {\hspace{-5.8cm}$X \enspace mg/l \enspace (Unknown)$} node [below] {\hspace{-5.8cm}$100 \enspace mg/l$};
    \path [arrow] (Process) -- (dummy_out)  node [above] {\hspace{-4cm}$A \enspace mg/l \enspace (Given)$} node [below] {\hspace{-3.9cm}($100-RE\%)\enspace mg/l$};
   \draw[arrow] (Process) -- (Removal);
\end{tikzpicture}
\end{figure}
Using the fact that if the inlet concentration was 100 mg/l, the outlet concentration would be 100 minus the removal efficiency.\\
Setup the equation as:  $\frac{In}{Out}: \enspace \frac{X \enspace mg/l}{A \enspace mg/l}=\frac{100}{100-RE\%}$\\
\vspace{0.3cm}
Calculate X using cross multiplication - if $\frac{A}{B}=\frac{C}{D} \implies A=B*\frac{C}{D}$:\\
$X \enspace mg/l=A \enspace mg/l*\frac{100}{100-RE\%}$\\

\pagebreak

\begin{enumerate}



\item  Given the following for a primary sedimentation tank: TSS removal efficiency = 63\% Effluent TSS concentration = 95 mg/L Calculate the influent TSS concentration (mg/L) Correct Answer(s): \\

*a. 257 \\
b. 356 \\
c. 94 \\
d. 220 \\

Solution:\\

\tikzstyle{block} = [rectangle, draw, fill=red!40, 
    text width=7em, text centered, rounded corners, minimum height=3em]
\tikzstyle{arrow} = [draw, -latex']
\begin{figure}[!h]
\centering
\begin{tikzpicture}[node distance =1.5cm, auto]
    \draw ++(0,0) node [block] (Process) {Primary \\Sedimentation\\ Tank};
   \node[node distance=1.5in] (dummy_in) [left of=Process] {In TSS};
   \node[node distance=1.5in] (dummy_out) [right of=Process] {Out TSS};
	\node (Removal) [below of=Process, yshift=-0.25in] {$Removal \enspace Efficiency=63\%$};
    \path [arrow] (dummy_in)-- (Process)  node [above] {\hspace{-4.39cm}$Xmg/l$} node [below] {\hspace{-4.39cm}$100mg/l$};
    \path [arrow] (Process) -- (dummy_out)  node [above] {\hspace{-3.cm}95mg/l} node [below] {\hspace{-3cm}37mg/l};
   \draw[arrow] (Process) -- (Removal);
\end{tikzpicture}
\end{figure}

$\frac{In}{Out} \enspace : \enspace \frac{Actual \enspace inlet \enspace  (X)}{95}=\frac{100}{100-63}$\\
$\implies \frac{Actual \enspace inlet \enspace  (X)}{95}=2.7$\\    
Rearranging the equation:   $Actual \enspace inlet (X)=2.7*95 = \boxed{257 mg/l}$\\



\item What is the clarifier influent TSS if its outlet concentration is 60 mg/l and the known clarifier removal efficiency is 75\%? \\
*a. 240 mg/l \\
b. 300 mg/l \\
c. 80 mg/l \\
d. 120 mg/l \\

Solution:\\
\tikzstyle{block} = [rectangle, draw, fill=red!40, 
    text width=7em, text centered, rounded corners, minimum height=3em]
\tikzstyle{arrow} = [draw, -latex']
\begin{figure}[!h]
\centering
\begin{tikzpicture}[node distance =1.5cm, auto]
    \draw ++(0,0) node [block] (Process) {Clarifier};
   \node[node distance=1.5in] (dummy_in) [left of=Process] {In TSS};
   \node[node distance=1.5in] (dummy_out) [right of=Process] {Out TSS};
	\node (Removal) [below of=Process, yshift=-0in] {$Removal \enspace Efficiency=75\%$};
    \path [arrow] (dummy_in)-- (Process)  node [above] {\hspace{-4.39cm}$Xmg/l$} node [below] {\hspace{-4.39cm}$100mg/l$};
    \path [arrow] (Process) -- (dummy_out)  node [above] {\hspace{-3.cm}60mg/l} node [below] {\hspace{-3cm}37mg/l};
   \draw[arrow] (Process) -- (Removal);
\end{tikzpicture}
\end{figure}

$\frac{In}{Out} \enspace : \enspace \frac{Actual \enspace inlet \enspace  (X)}{60}=\frac{100}{100-75}$\\
$\implies \frac{Actual \enspace inlet \enspace  (X)}{60}=4$\\    
Rearranging the equation:   $Actual \enspace inlet (X)=4*60 = \boxed{240 mg/l}$\\



\item If a primary clarifier consistently operates at 30\% efficiency and produces an effluent which averages 140 mg/l BOD, what is the influent BOD?\\
Solution:\\
\vspace{0.3cm}
\tikzstyle{block} = [rectangle, draw, fill=red!40, 
    text width=7em, text centered, rounded corners, minimum height=3em]
\tikzstyle{arrow} = [draw, -latex']
\begin{figure}[!h]
\centering
\begin{tikzpicture}[node distance =1.5cm, auto]
    \draw ++(0,0) node [block] (Process) {Clarifier};
   \node[node distance=1.5in] (dummy_in) [left of=Process] {In TSS};
   \node[node distance=1.5in] (dummy_out) [right of=Process] {Out TSS};
	\node (Removal) [below of=Process, yshift=-0in] {$Removal \enspace Efficiency=30\%$};
    \path [arrow] (dummy_in)-- (Process)  node [above] {\hspace{-4.39cm}$Xmg/l$} node [below] {\hspace{-4.39cm}$100mg/l$};
    \path [arrow] (Process) -- (dummy_out)  node [above] {\hspace{-3.cm}140mg/l} node [below] {\hspace{-3cm}70mg/l};
   \draw[arrow] (Process) -- (Removal);
\end{tikzpicture}
\end{figure}

$\frac{In}{Out} \enspace : \enspace \frac{Actual \enspace inlet \enspace  (X)}{140}=\frac{100}{100-30}$\\
$\implies Actual \enspace inlet \enspace  (X)=\frac{140*100}{70}=\boxed{200mg/l}$\\

\item What is the clarifier removal efficiency if the inlet and outlet concentrations are 300 mg/l and 60 mg/l respectively?\\

a. 20\%\\
b. 30\%\\
*c. 80\%\\
d. 72\%


\item A primary clarifier has an influent suspended solids concentration of 250 mg/L. If the suspended solids removal efficiency is 60\%, what is the primary effluent suspended solids concentration?\\

a. 10 mg/L\\
*b. 100 mg/L\\
c. 50 mg/L\\
d. 150 mg/L

\item The influent to a trickling filter plant is 200 mg/L and the effluent BOD is 20 mg/L. What is the BOD removal efficiency (\%)?
Correct Answer(s):\\
a. 90.0


\item If a plant removes 35\% of the influent BOD in the primary treatment and 85\% of the remaining BOD in the secondary system, what is the BOD removal (lbs/day), if the BOD of the final effluent is 20mg/l 

a. 2445 lbs/day 
*b. 2211 lbs/day 
c. 3800 lbs/day 
d. 4200 lbs/day 


\item If a plant removes 35\% of the influent BOD in the primary treatment and 85\% of the remaining BOD in the secondary system, what is the BOD of the raw wastewater if the BOD of the final effluent is 20mg/l \\

a. 48 mg/l \\
b. 67 mg/l \\
c. 86.7 mg/l \\
d. 155 mg/l \\
*e. 205 mg/l
Solution:\\

\begin{figure}[!h]
\centering
\begin{tikzpicture}[node distance =1.5cm, auto]
    \draw ++(0,0) node [block] (Primary) {Primary};
    
   \node[node distance=1.9in] (dummy_in) [left of=Primary] {Influent BOD};
   \node[node distance=1.9in] (dummy_out) [right of=Primary] {Primary BOD Out};
	\node (Removal) [below of=Primary, yshift=-0in] {$Removal \enspace Efficiency=35\% $};
    \path [arrow] (dummy_in)-- (Primary)  node [above] {\hspace{-4.8cm}$X \enspace mg/l \enspace$} node [below] {};
    \path [arrow] (Primary) -- (dummy_out)  node [above] {\hspace{-4.9cm}$0.65X \enspace mg/l$} node [below] {};
   \draw[arrow] (Process) -- (Removal);
\end{tikzpicture}
\end{figure}


\begin{figure}[!h]
\centering
\begin{tikzpicture}[node distance =1.5cm, auto]
    \draw ++(0,0) node [block] (Secondary) {Secondary};
    
   \node[node distance=1.9in] (dummy_in) [left of=Secondary] {Primary BOD Out};
   \node[node distance=1.9in] (dummy_out) [right of=Secondary] {Secondary BOD Out};
	\node (Removal) [below of=Secondary, yshift=-0in] {$Removal \enspace Efficiency=85\% $};
    \path [arrow] (dummy_in)-- (Secondary)  node [above] {\hspace{-4.8cm}$0.65X \enspace mg/l \enspace$} node [below] {\hspace{-5cm}$100 \enspace mg/l$};
    \path [arrow] (Secondary) -- (dummy_out)  node [above] {\hspace{-4.9cm}$20 \enspace mg/l$} node [below] {\hspace{-4.9cm}$15 \enspace mg/l$};
   \draw[arrow] (Process) -- (Removal);
\end{tikzpicture}
\end{figure}
\vspace{0.3cm}
For the Secondary process:\\
$\frac{In}{Out}: \enspace \frac{0.65X}{20}=\frac{100}{15} \implies X \enspace mg/l=\frac{100*20}{15*0.65}=\boxed{205 \enspace mg/l}$\\

\vspace{0.3cm}
Alternate Solution \#1

$\xrightarrow[
				\text{X}\frac{mg}{l}
			]
			{
			\text{Influent BOD}
			}
 \boxed{Primary}
 \xrightarrow[
 				\text{X-0.35X=X*(1-0.35)=0.65X}\frac{mg}{l}
 			]
 			{
 			\text{Primary Effluent BOD}
 			}
 \boxed{Secondary}
 \xrightarrow[
				\text{0.65X-0.5525X=(0.65-0.5525)X=0.0975X }
			 ]
			{
			\text{Secondary Effluent BOD}
			}
$\\

\hspace{2.8cm}$\downarrow$ {\tiny(0.35X)BOD Removed}\hspace{3.2cm}$\downarrow$ {\tiny(0.65*0.85)X = 0.5525X BOD Removed}\\
$\implies 0.0975X=20 \implies X=\frac{20}{0.0975}=\boxed{205\frac{mg}{l}}$\\

\vspace{0.3cm}

Alternate Solution \#2:\\
$\xrightarrow[\text{X}\frac{mg}{l}]{\text{Influent BOD}}\boxed{Primary}\xrightarrow[\text{0.65X}]{\text{Primary Effluent BOD}}\boxed{Secondary}\xrightarrow[\text{(0.65*0.15)X}]{\text{Secondary Effluent BOD}}$\\
\hspace{2.8cm}$\downarrow$ {\tiny(0.35X)BOD Removed}\hspace{2.2cm}$\downarrow$ {\tiny(0.65X*0.85)BOD Removed}\\

Primary Effluent BOD = Influent BOD * (1-Primary BOD Removal), and\\
Secondary Effluent BOD=[Primary Effluent BOD]*(1-Secondary BOD Removal)\\
Secondary Eff. BOD=[Influent BOD * (1-Primary BOD Removal)]*(1-Secondary BOD Removal)\\

Therefore, 20 = [X*(1-0.35)] * (1-0.85)= X*0.65*0.15\\
$\implies 20 \enspace \frac{mg}{l}= 0.0975X \implies X=\frac{20}{0.0975}=\boxed{205 \enspace \frac{mg}{l}}$\\

\item Calculate the inlet concentration if the outlet concentration is 80 mg/l and the process removal efficiency is 60\%\\

\tikzstyle{block} = [rectangle, draw, fill=red!40, 
    text width=6em, text centered, rounded corners, minimum height=3em]
\tikzstyle{arrow} = [draw, -latex']
\begin{figure}[!h]
\centering
\begin{tikzpicture}[node distance =1.5cm, auto]
    \draw ++(0,0) node [block] (Process) {Process};
   \node[node distance=1.5in] (dummy_in) [left of=Process] {In};
   \node[node distance=1.5in] (dummy_out) [right of=Process] {Out};
	\node (Removal) [below of=Process, yshift=-0in] {$Removal \enspace Efficiency=60\%$};
    \path [arrow] (dummy_in)-- (Process)  node [above] {\hspace{-4.39cm}$Xmg/l$} node [below] {\hspace{-4.39cm}$100mg/l$};
    \path [arrow] (Process) -- (dummy_out)  node [above] {\hspace{-3.cm}80mg/l} node [below] {\hspace{-3cm}40mg/l};
   \draw[arrow] (Process) -- (Removal);
\end{tikzpicture}
\end{figure}

$\frac{In}{Out} \enspace : \enspace \frac{Actual \enspace inlet \enspace  (X)}{80}=\frac{100}{100-60}$\\
$\implies \frac{Actual \enspace inlet \enspace  (X)}{80}=2.5$\\    
Rearranging the equation:   $Actual \enspace inlet (X)=2.5*80 = \boxed{200 mg/l}$\\

\item Calculate the outlet concentration if the inlet concentration is 80 mg/l and the process removal efficiency is 60\%\\
Solution:\\

\tikzstyle{block} = [rectangle, draw, fill=red!40, 
    text width=6em, text centered, rounded corners, minimum height=3em]
\tikzstyle{arrow} = [draw, -latex']
\begin{figure}[!h]
\centering
\begin{tikzpicture}[node distance =1.5cm, auto]
    \draw ++(0,0) node [block] (Process) {Process};
   \node[node distance=1.5in] (dummy_in) [left of=Process] {In};
   \node[node distance=1.5in] (dummy_out) [right of=Process] {Out};
	\node (Removal) [below of=Process, yshift=-0in] {$Removal \enspace Efficiency=60\%$};
    \path [arrow] (dummy_in)-- (Process)  node [above] {\hspace{-4.39cm}$80mg/l$} node [below] {\hspace{-4.39cm}$100mg/l$};
    \path [arrow] (Process) -- (dummy_out)  node [above] {\hspace{-3.cm}$Xmg/l$} node [below] {\hspace{-3cm}40mg/l};
   \draw[arrow] (Process) -- (Removal);
\end{tikzpicture}
%\caption[MFCC]{Diagrama en bloques del cálculo de las MFCC para un frame.}
%\label{MFCC}
\end{figure}
\vspace{0cm}
$\frac{Out}{In} \enspace:\enspace\frac{Actual \enspace Outlet (X)}{80}=\frac{100-60}{100}$\\
$\implies \frac{Actual \enspace Outlet (X)}{80} =0.4$\\
$\implies Actual \enspace  Outlet (X) = 0.4 * 80 = \boxed{32 mg/l}$\\

\item Calculate the primary clarifier influent solids concentration if its outlet concentration is 60 mg/l and the known clarifier removal efficiency is 75\%?\\
$\frac{Actual \enspace  inlet \enspace (X)}{Actual \enspace outlet}=\frac{100}{100-Removal \enspace efficiency}$\\ 
$\frac{Actual \enspace  inlet \enspace (X)}{60}=\frac{100}{100-75}=4$\\
$\implies Actual \enspace inlet \enspace (X)=4*60 = \boxed{240 mg/l}$\\

\item What is the \% removal efficiency if the influent concentration is 10 mg/L and the effluent concentration is 2.5 mg/L?\\
$Removal \enspace Rate (\%) = \frac{In-Out}{In}*100 \implies \frac{10-2.5}{10}*100=\boxed{75\%}$

\item The influent to a trickling filter plant is 200 mg/L and the effluent BOD is 20 mg/L. What is the BOD removal efficiency (\%)?\\
Solution:\\
\vspace{0.3cm}
\tikzstyle{block} = [rectangle, draw, fill=red!40, 
    text width=6em, text centered, rounded corners, minimum height=3em]
\tikzstyle{arrow} = [draw, -latex']
\begin{figure}[!h]
\centering
\begin{tikzpicture}[node distance =1.5cm, auto]
    \draw ++(0,0) node [block] (Process) {Process};
   \node[node distance=1.5in] (dummy_in) [left of=Process] {};
   \node[node distance=1.5in] (dummy_out) [right of=Process] {};
	\node (Removal) [below of=Process, yshift=-0in] {$Removal \enspace Efficiency=?\%$};
    \path [arrow] (dummy_in)-- (Process)  node [above] {\hspace{-4.39cm}$200mg/l$} node [below] {};
    \path [arrow] (Process) -- (dummy_out)  node [above] {\hspace{-1.8cm}$20mg/l$} node [below] {};
   \draw[arrow] (Process) -- (Removal);
\end{tikzpicture}
\end{figure}
\vspace{0.5cm}
$Removal \enspace Efficiency \enspace (\%) = \frac{In-Out}{In}*100 \implies \frac{200-20}{200}*100=\boxed{90\%}$




\item Calculate the primary clarifier influent solids concentration if its outlet concentration is 60 mg/l and the known clarifier removal efficiency is 75\%?\\
$\frac{Actual \enspace  inlet \enspace (X)}{Actual \enspace outlet}=\frac{100}{100-Removal \enspace efficiency}$\\ 
$\frac{Actual \enspace  inlet \enspace (X)}{60}=\frac{100}{100-75}=4$\\
$\implies Actual \enspace inlet \enspace (X)=4*60 = \boxed{240 mg/l}$\\

\item If a primary clarifier consistently operates at 30\% efficiency and produces an effluent which averages 140 mg/l BOD, what is the influent BOD? \\
Solution:\\
\vspace{0.3cm}
\tikzstyle{block} = [rectangle, draw, fill=red!40, 
    text width=6em, text centered, rounded corners, minimum height=3em]
\tikzstyle{arrow} = [draw, -latex']
\begin{figure}[!h]
\centering
\begin{tikzpicture}[node distance =1.5cm, auto]
    \draw ++(0,0) node [block] (Process) {Process};
   \node[node distance=1.5in] (dummy_in) [left of=Process] {In};
   \node[node distance=1.5in] (dummy_out) [right of=Process] {Out};
	\node (Removal) [below of=Process, yshift=-0in] {$Removal \enspace Efficiency=30\%$};
    \path [arrow] (dummy_in)-- (Process)  node [above] {\hspace{-4.39cm}$Xmg/l$} node [below] {\hspace{-4.39cm}$100mg/l$};
    \path [arrow] (Process) -- (dummy_out)  node [above] {\hspace{-3.cm}$140mg/l$} node [below] {\hspace{-3cm}70mg/l};
   \draw[arrow] (Process) -- (Removal);
\end{tikzpicture}
%\caption[MFCC]{Diagrama en bloques del cálculo de las MFCC para un frame.}
%\label{MFCC}
\end{figure}
$\frac{Actual \enspace  inlet \enspace (X)}{Actual \enspace outlet}=\frac{100}{100-Removal \enspace efficiency}$\\ 
$\frac{Actual \enspace  inlet \enspace (X)}{140}=\frac{100}{100-30}=1.43$\\
$\implies Actual \enspace inlet \enspace (X)=1.43*140 = \boxed{200 mg/l}$\\



\item If a primary clarifier consistently operates at 30 \% efficiency and produces an effluent which averages 140 mg/l BOD, what is the influent BOD? 

a. 100 mg/l \\
b. 125 mg/l \\
c. 175 mg/l \\
*d. 200 mg/l \\
e. 225 mg/l \\

\item A primary clarifier has an influent suspended solids concentration of 250 mg/L. If the suspended solids removal efficiency is 60\%, what is the primary effluent suspended solids concentration \\

a. 10 mg/L \\
*b. 100 mg/L \\
c. 50 mg/L \\
d. 150 mg/L \\

\item The influent to a trickling filter plant is 200 mg/L and the effluent BOD is 20 mg/L. What is the BOD removal efficiency (\%)? \\

a. 20\% \\
*b. 90\% \\
c. 10\% \\
d. 180\% \\

\item A 3.1 MGD flow with a 190 mg/l TSS concentration is treated in a primary clarifier which averages 55\% removal efficiency. Calculate the pounds/day TSS in the primary effluent. \\

a. 2701 lbs/day \\
*b. 2211 lbs/day \\
c. 3800 lbs/day \\
d. 4200 lbs/day \\

\item Calculate the outlet concentration if the inlet concentration is 80 mg/l and the process removal efficiency is 60\% \\

*a. 32mg/l \\
b. 63mg/l \\
c. 20mg/l \\
d. 48mg/l \\

\item What is the \% removal efficiency if the influent concentration is 10 mg/L and the effluent concentration is 2.5 mg/L? \\

a. 25\% \\
b. 97.5\% \\
*c. 75\% \\
d. 2.5\% \\

\item A rectangular sedimentation tank is 85 feet long, 35 feet wide, and 14 feet deep including 3 feet of freeboard. Flow to this tank is 2.3 MGD. Calculate the surface loading to this tank in gpd per ft2. \\

a. 318 gpd/ft2 \\
*b. 773 gpd/ft2 \\
c. 845 gpd/ft2 \\
d. 1932 gpd/ft2 \\

\item  If a primary clarifier consistently operates at 30 \% efficiency and produces an effluent which averages 140 mg/l BOD, what is the influent BOD? \\

a. 100 mg/l \\
b. 125 mg/l \\
c. 175 mg/l \\
*d. 200 mg/l \\

\item  A primary clarifier has an influent suspended solids concentration of 250 mg/L. If the suspended solids removal efficiency is 60\%, what is the primary effluent suspended solids concentration \\

a. 10 mg/L \\
*b. 100 mg/L \\
c. 50 mg/L \\
d. 150 mg/L \\

\item  The influent to a trickling filter plant is 200 mg/L and the effluent BOD is 20 mg/L. What is the BOD removal efficiency (\%)? \\

a. 20\% \\
*b. 90\% \\
c. 10\% \\
d. 180\% \\

\item A 3.1 MGD flow with a 190 mg/l TSS concentration is treated in a primary clarifier which averages 55\% removal efficiency. Calculate the pounds/day TSS in the primary effluent. \\

a. 2701 lbs/day \\
*b. 2211 lbs/day \\
c. 3800 lbs/day \\
d. 4200 lbs/day \\



\item Calculate the outlet concentration if the inlet concentration is 80 mg/l and the process removal efficiency is 60\% \\

*a. 32mg/l \\
b. 63mg/l \\
c. 20mg/l \\
d. 48mg/l \\

\item What is the \% removal efficiency if the influent concentration is 10 mg/L and the effluent concentration is 2.5 mg/L? \\

a. 25\% \\
b. 97.5\% \\
*c. 75\% \\
d. 2.5\% \\



\item Use the information below to answer the following questions.\\
Plant Flow: 4.5 MGD\\
Influent P concentration: 1.5 mg/l\\
Aeration tank volume: 2000 cu.ft.\\
Effluent P concentration: 0.5 mg/l\\
What is the phosphorus removal efficiency?:\\
a. 34 percent\\
b. 67 percent\\
c. 1.0 percent\\
d. 10 percent\\
Solution:\\
$\frac{1.5-0.5}{1.5}*100=\boxed{67\%}$

\item Calculate the log removal for a water treatment plant if the samples show a raw water coliform count of 295/100 mL (through extrapolation) and the finished water shows 2.0/100 mL.
a. 1.8 log removal
b. 2 log removal
c. 2.1 log removal
*d. 2.2 log removal

\tikzstyle{block} = [rectangle, draw, fill=red!40, 
    text width=6em, text centered, rounded corners, minimum height=3em]
\tikzstyle{arrow} = [draw, -latex']
\begin{figure}[!h]
\centering
\begin{tikzpicture}[node distance =1.5cm, auto]
    \draw ++(0,0) node [block] (Process) {Treatment};
   \node[node distance=1.9in] (dummy_in) [left of=Process] {295/100 mL};
   \node[node distance=1.9in] (dummy_out) [right of=Process] {2.0/100 mL};
	\node (Removal) [below of=Process, yshift=-0in] {$Coliform \enspace Removal - log \enspace Reduction $};
    \path [arrow] (dummy_in)-- (Process)  node [above] {} node [below] {};
    \path [arrow] (Process) -- (dummy_out)  node [above] {} node [below] {};
   \draw[arrow] (Process) -- (Removal);
\end{tikzpicture}
\end{figure}


Log Reduction = log$_{10}$C$_{In}$ \hspace{0.2cm} - \hspace{0.2cm}log$_{10}$ C$_{Out}$\hspace{0.2cm}= \hspace{0.2cm}log$_{10}$ \hspace{0.02cm} 295 - log$_{10}$\hspace{0.02cm} 2\hspace{0.2cm}=\hspace{0.2cm}2.47\hspace{0.2cm}-\hspace{0.2cm}0.30\hspace{0.2cm}=\hspace{0.2cm}$\boxed{2.17 \enspace log \enspace removal}$

  \item Which is the percentage of removal across a settling basin, if the influent is 17.1 NTU and the effluent is 1.13 NTU ?\\
  


  \item Calculate the percent removal across a settling basin and filter complex, if the raw water influent is $5.45 \mathrm{ntu}$ and the effluent (post filters) is $0.018 \mathrm{ntu}$. Give answer to three significant figures.\\
  
      \item Calculate the amount of iron removed in pounds per year from a water plant that treats an average of $20.2 \mathrm{mgd}$ if the average iron concentration is 0.52 mg/l and the removal efficiency is 84 \%.\\
a. $26,859 \mathrm{lb} / \mathrm{yr}$ of Fe removed\\
*b. $27,000 \mathrm{lb} / \mathrm{yr}$ of Fe removed\\
c. $31,975 \mathrm{lb} / \mathrm{yr}$ of Fe removed\\
d. $32,000 \mathrm{lb} / \mathrm{yr}$ of Fe removed\\

\begin{figure}[]
\begin{tikzpicture}
    \newcommand{\R}{1.5}

\path[help lines,step=.2] (0,0) grid (16,3);
\path[help lines,line width=.6pt,step=1] (0,0) grid (16,3);
%\foreach \x in {0,1,2,3,4,5,6,7,8,9,10,11,12,13,14,15,16}
%\node[anchor=north] at (\x,0) {\x};
%\foreach \y in {0,1,2,3,4,5,6}
%\node[anchor=east] at (0,\y) {\y};
%-------------CIRCLE-----------------------------------
\draw[black,fill=gray!10] (8,3) circle (\R);
\draw[black, very thick, rotate=0](6.5,3) -- (9.5,3);
\draw (8,3.6) node[text width=3cm,align=center]
  {\scriptsize{lbs/yr}};
\draw (7.1,2.5) node[text width=3cm,align=center]
  {\tiny{(0.52*0.84)\\ mg/l}};
\draw (8.9,2.5) node[text width=3cm,align=center]
  {\tiny{20.2 MGD}};
  \draw (8,2)node[text width=3cm,align=center]
  {\tiny{8.34}};
\draw[black, very thick, rotate=0](7.2,1.7) -- (8,3);
\draw[black, very thick, rotate=0](8.8,1.7) -- (8,3);
\end{tikzpicture}
\end{figure}
$\dfrac{\mathrm{lbs}}{\mathrm{yr}}=\mathrm{Flow}\dfrac{{\mathrm{MG}}}{\mathrm{yr}}* \mathrm{Concentration \enspace iron \enspace removed}\dfrac{\mathrm{mg}}{\mathrm{l}}*8.34$
\\
\vspace{0.2cm}
$\dfrac{\mathrm{lbs}}{\mathrm{day}}=20.2 \dfrac{\mathrm{MG}}{\cancel{\mathrm{day}}}* 365\dfrac{\cancel{\mathrm{days}}}{\mathrm{yr}}*(20.2*0.84)\dfrac{\mathrm{mg}}{\mathrm{l}}*8.34 = \boxed{26,859\dfrac{lbs}{yr}}$

   \item Determine the amount of iron removed per year, if the iron concentration is 0.21 mg/l, the plant treats an average of $14.1 \mathrm{mgd}$, and the removal efficiency is $95.7 \%$ (0.957).\\
a. $8,000 \mathrm{lb} / \mathrm{yr}$\\
b. $8,200 \mathrm{lb} / \mathrm{yr}$\\
*c. $8,600 \mathrm{lb} / \mathrm{yr}$\\
d. $9,000 \mathrm{lb} / \mathrm{yr}$\\

A: $(0.21*0.957)\dfrac{mg}{l}*14.1\dfrac{MG}{day}*8.34*365\dfrac{days}{yr}=\boxed{8,626\enspace lbs/day}$\\
 

\end{enumerate}


\end{document}